% Source: https://artofproblemsolving.com/community/c3h2020601_perimeter_of_monster
% Stable identity: aops-topic-2020601-q01
The boundary of the monster consists of two radii and the major arc. The arc subtends $360^\circ-60^\circ=300^\circ$ at the centre $O$.

\begin{asy}
usepackage("amsmath");
size(4.5cm);
defaultpen(fontsize(10pt));
pair O=(0,0);
filldraw(O--arc(O,1,30,330)--cycle,gray(0.85));
draw(arc(O,1,30,330),red+linewidth(1.3));
draw(O--dir(30),blue+linewidth(1.3)); draw(O--dir(330),blue+linewidth(1.3));
draw(arc(O,0.3,-30,30)); label("$60^\circ$",(0.55,0));
label("$300^\circ$",(-0.45,0.12));
dot(O); label("$O$",O,2*S);
label("$1$",dir(30)/2,NW,blue); label("$1$",dir(330)/2,SW,blue);
label("arc $=\dfrac{5\pi}{3}$",dir(200),W,red);
label("\mbox{}",(min(currentpicture,user=true)+max(currentpicture,user=true))/2,black); // keeps the text after the figure black when LaTeX typesets the labels inline
\end{asy}

The two radii contribute $1+1=2$ cm. The arc is $\dfrac{300}{360}=\dfrac{5}{6}$ of the circumference $2\pi$, i.e. $\dfrac{5\pi}{3}$ cm. The perimeter is
\[
\boxed{2+\dfrac{5\pi}{3}}\text{ cm}\approx7.24\text{ cm}.
\]
